% Substitua o texto abaixo pelo seu resumo em inglês.

This work presents the design and implementation of an Embedded Systems infrastructure for the centralized monitoring of experiments in Physics teaching laboratories. The motivation is based on the structural barriers of traditional laboratories, which often operate in an isolated manner and require dedicated hardware for each workbench, hindering pedagogical management and scalability. The proposed architecture uses ESP32 microcontrollers as edge nodes for data acquisition and a Raspberry Pi server as a central coordinator. The methodology followed a technological development approach, evaluated through a Proof of Concept (PoC) with the Capacitor Charge and Discharge experiment (RC Circuit). Communication between devices occurs via HTTP/JSON protocol over a Wi-Fi network, eliminating the need for wired connections and centralizing telemetry in a PostgreSQL database. The system features a web interface developed in Python (Quart), allowing real-time visualization and data download by students and teachers. The results demonstrate the feasibility of an orchestration layer that decouples the user interface from the acquisition hardware, promoting an integrated, scalable laboratory environment capable of supporting hybrid and remote teaching practices.

\textbf{Key-words}: Embedded Systems; Wireless Experimentation; Physics Teaching; ESP32; Telemetry.
